
\documentclass{article}
\usepackage{colortbl}
\usepackage{makecell}
\usepackage{multirow}
\usepackage{supertabular}

\begin{document}

\newcounter{utterance}

\centering \large Interaction Transcript for game `matchit\_ascii', experiment `similar\_grid\_2', episode 2 with Qwen3{-}30B{-}A3B{-}t0.0.
\vspace{24pt}

{ \footnotesize  \setcounter{utterance}{1}
\setlength{\tabcolsep}{0pt}
\begin{supertabular}{c@{$\;$}|p{.15\linewidth}@{}p{.15\linewidth}p{.15\linewidth}p{.15\linewidth}p{.15\linewidth}p{.15\linewidth}}
    \# & $\;$A & \multicolumn{4}{c}{Game Master} & $\;\:$B\\
    \hline

    \theutterance \stepcounter{utterance}  
    & & \multicolumn{4}{p{0.6\linewidth}}{
        \cellcolor[rgb]{0.9,0.9,0.9}{
            \makecell[{{p{\linewidth}}}]{
                \texttt{\tiny{[P1$\langle$GM]}}
                \texttt{You are participating in a collaborative guessing game.} \\
\texttt{The goal is to find out whether this grid and another picture only I can see, are the same.} \\
\texttt{Please describe your grid first. Then, I will provide my description and we can ask each other questions about the grids to figure out whether they are the same. Now start your short grid description with "DESCRIPTION:" followed by the description. Do not add anything else. This is your grid:} \\
\texttt{X X ▢ X X} \\
\texttt{X X ▢ X X} \\
\texttt{X X X X X} \\
\texttt{X X ▢ X X} \\
\texttt{X X ▢ X X} \\
            }
        }
    }
    & \\ \\

    \theutterance \stepcounter{utterance}  
    & \multicolumn{4}{p{0.6\linewidth}}{
        \cellcolor[rgb]{0.9,0.9,0.9}{
            \makecell[{{p{\linewidth}}}]{
                \texttt{\tiny{[P1$\rangle$GM]}}
                \texttt{<think>} \\
\texttt{Okay, let's see. The user wants me to describe my grid first. The grid they provided is:} \\
\\ 
\texttt{X X ▢ X X} \\
\texttt{X X ▢ X X} \\
\texttt{X X X X X} \\
\texttt{X X ▢ X X} \\
\texttt{X X ▢ X X} \\
\\ 
\texttt{So, I need to make sure I accurately represent this grid. Let me count the rows and columns. There are 5 rows and 5 columns. Each row has 5 elements. The third row is all X's, which is different from the others. The other rows have three X's and two empty squares (represented by ▢). Wait, no, looking again: each of the first, second, fourth, and fifth rows has X X ▢ X X. So that's two X's, then an empty, then two X's. The third row is all X's. So the grid has a horizontal line of X's in the middle. } \\
\\ 
\texttt{I need to describe this grid in a concise way. The user said to start with "DESCRIPTION:" followed by the description. So I should just write that. Let me check if there's any ambiguity. Maybe mention the number of rows and columns, and the pattern. For example, "A 5x5 grid with the third row completely filled with X's, and the other rows having X's in the first, second, fourth, and fifth positions, with an empty space in the third position." But maybe the user} \\
            }
        }
    }
    & & \\ \\

    \theutterance \stepcounter{utterance}  
    & & & \multicolumn{2}{p{0.3\linewidth}}{
        \cellcolor[rgb]{0.9,0.9,0.9}{
            \makecell[{{p{\linewidth}}}]{
                \texttt{\tiny{[GM$|$GM]}}
                \texttt{abort, first word: <think>, but expected DESCRIPTION:} \\
            }
        }
    }
    & & \\ \\

    \theutterance \stepcounter{utterance}  
    & & & \multicolumn{2}{p{0.3\linewidth}}{
        \cellcolor[rgb]{0.9,0.9,0.9}{
            \makecell[{{p{\linewidth}}}]{
                \texttt{\tiny{[GM$|$GM]}}
                \texttt{Aborted} \\
            }
        }
    }
    & & \\ \\

\end{supertabular}
}

\end{document}
